\documentclass[11pt]{letter}
\usepackage[T1]{fontenc}
\usepackage[utf8]{inputenc}
\usepackage[margin=2.5cm]{geometry}
\usepackage[hidelinks]{hyperref}
\signature{Uday Pratap Singh\\(on behalf of all authors)}
\address{Department of Artificial Intelligence and Data Science\\Poornima Institute of Engineering and Technology\\Jaipur 302017, Rajasthan, India\\udaypratap@allduniv.ac.in}
\begin{document}
\begin{letter}{The Editor-in-Chief\\Computer Networks}
\opening{Dear Editor,}

We submit our manuscript ``Choosing backup paths from hazard history: disaster-aware path selection for traffic-engineering recovery, tested on real cyclones and earthquakes'' for consideration as a full research article in \emph{Computer Networks}.

The paper deals with network reliability and routing, topics 5 and 2 of the journal's scope. Traffic-engineering controllers recover from failures by moving traffic onto tunnels they computed in advance; when a cyclone or earthquake destroys all of a pair's tunnels at once, that traffic is lost although the network may still connect the pair. We ask whether the public record of past hazards can tell an operator which tunnels to precompute, and we answer it with real data: USGS ShakeMap intensity grids, NOAA HURDAT2 and IBTrACS cyclone tracks with wind radii, and seven real backbone topologies.

The proposed rule, DART, keeps link-disjoint tunnels and replaces the last one only when past footprints show a measurable benefit. It was developed on two networks, frozen, and tested twice under criteria committed to a public repository before each test. On 211 real cyclone and earthquake footprints from 2018--2025 it lost 0.16\% of demand on average against 0.30\% for edge-disjoint tunnels (Holm-corrected $p=8.6\times10^{-8}$), at no cost in normal-operation utilisation and with no additional online computation.

We have tried to report the limits as clearly as the gains. The absolute improvement is small compared with the traffic that no routing can recover after a disaster; real earthquakes broke too few backbone links to separate any of the rules; and the first version of the method failed its pilot, which we document. All code, data, configuration files and the commit history of the pre-registered criteria are public at \url{https://github.com/singhudaypratap/dart-disaster-aware-te}, and every number in the paper is regenerated by one script from committed result files.

This manuscript is original, has not been published, and is not under consideration elsewhere. All authors have approved the submission and declare no competing interests. An AI assistant (Claude, Anthropic) was used for code development, analysis and drafting, as stated in the manuscript's declaration; the authors checked all content and take full responsibility for it.

Thank you for considering our work.

\closing{Sincerely,}
\end{letter}
\end{document}
